\section{Limits of Outcome-Only and Prefix-Local Supervision}

POV-OPD is motivated by two complementary limitations. A terminal outcome is reliable at the trajectory level but does not identify token-specific credit within one sampled response. Conversely, a teacher supplies dense token-level preferences, but global teacher superiority does not imply that every local preference remains reliable on student-generated prefixes. We formalize these two points below.

\subsection{A Terminal Outcome Does Not Identify Token Credit}

Let $Y=(y_1,\ldots,y_T)$ be a response sampled from the student policy, and let $R(Y)\in\{0,1\}$ be its verifier outcome. The ideal outcome credit of action $y_t$ at prefix $h_t=(x,y_{<t})$ is its conditional advantage
\begin{equation}
A_t^R = Q^{\pi}(h_t,y_t)-V^{\pi}(h_t),
\label{eq:ideal-token-credit}
\end{equation}
where $Q^{\pi}(h_t,y_t)$ is the expected final outcome after taking $y_t$ and following $\pi$. Outcome-only methods instead construct a Monte Carlo estimate from the single terminal observation, typically
\begin{equation}
\widehat A_t^{\mathrm{out}}=R(Y)-b(x), \qquad t=1,\ldots,T,
\label{eq:outcome-broadcast}
\end{equation}
and broadcast the same scalar to every response token.

\paragraph{Proposition 1.}
For $T>1$, the observed pair $(Y,R(Y))$ is insufficient to identify the token-level credits $\{A_t^R\}_{t=1}^{T}$ for an individual rollout.

\paragraph{Proof.}
Fix an observed rollout $Y$ and its terminal reward $R(Y)$. Consider any position $t$ and an alternative action $y_t'$. Two environments can assign exactly the same transitions and terminal reward along the observed rollout $Y$, while assigning different continuation rewards after the unobserved alternative $y_t'$. These environments therefore produce the same observation $(Y,R(Y))$, but induce different values of $Q^{\pi}(h_t,y_t')$, different baselines $V^{\pi}(h_t)$, and hence different action advantages in Equation~\eqref{eq:ideal-token-credit}. No estimator that observes only $(Y,R(Y))$ can distinguish the two environments, so it cannot recover the token-specific credits in both. \hfill $\square$

This proposition does not imply that outcome policy gradients are invalid: terminal Monte Carlo returns can be unbiased after averaging over sufficiently many trajectories. The limitation is finer-grained. Within one long response, the same outcome scalar cannot distinguish a useful reasoning step from a detour, a recoverable mistake, or a token whose effect is neutral. Outcome feedback is therefore reliable but sparse and high variance as a token-level learning signal.

\subsection{A Globally Strong Teacher Can Be Locally Unreliable}

Let $J(\pi;x)$ denote the success probability of policy $\pi$ on prompt $x$. OPD commonly assumes a teacher that is stronger on average:
\begin{equation}
\mathbb{E}_{x}[J(\pi_T;x)]>
\mathbb{E}_{x}[J(\pi_S;x)].
\label{eq:global-teacher-superiority}
\end{equation}
This aggregate condition does not imply superiority at every prompt or every prefix visited by the student.

\paragraph{Proposition 2.}
A teacher can satisfy Equation~\eqref{eq:global-teacher-superiority} while providing an outcome-conflicting local preference on a student-generated prefix.

\paragraph{Proof.}
Consider two prompt types. Prompt $x_1$ occurs with probability $0.9$, and prompt $x_2$ with probability $0.1$. Let the student's success probability be $0.5$ on both prompts. Let the teacher's success probabilities be $0.9$ on $x_1$ and $0.1$ on $x_2$. The teacher is globally stronger because
\begin{equation}
0.9\times0.9+0.1\times0.1=0.82>0.5.
\end{equation}
However, on $x_2$ the teacher is locally worse than the student. In a two-action realization, the teacher can assign higher probability to the incorrect action at a prefix of $x_2$. The corresponding teacher--student log-probability gap then favors an action that decreases conditional success. Thus global teacher superiority does not guarantee local outcome alignment on every student-visited prefix. \hfill $\square$

The issue becomes more pronounced for long reasoning traces. OPD queries the teacher on prefixes generated by the student, including prefixes that may be rare under the teacher. Although the teacher always defines a conditional distribution at such states, its local logit difference need not be a reliable direction for improving the final verifier outcome.

\subsection{Implication for POV-OPD}

The two propositions expose a complementary gap. Outcome-only learning supplies a trustworthy final label but cannot resolve token-level credit from a single rollout. Unvalidated OPD supplies dense local information but cannot guarantee that this information is outcome-aligned at every student-generated prefix. POV-OPD therefore does not replace one signal with the other. It uses final correctness to validate the direction of dense teacher preferences, and uses sustained teacher excess surprisal as a conservative indicator of prefix drift. The next section describes how these two validation signals calibrate the original dense OPD advantage.
